% Training. Sources: refs/docs/train_eval_data.md §1, refs/docs/model_arch.md §4; HTML §5.

\section{Training}
\label{sec:training}

AmphionASR~2.0 starts from Qwen3-ASR-1.7B and is trained in two stages, both updating all parameters except the frozen ones listed in \Cref{tab:training}.
Stage~1 consists of three rounds of multi-task supervised fine-tuning.
Stage~2 continues from the stage-1 model with target-speaker robustness training in five segments.
Each segment starts from the weights of the previous one, and its data recipe is adjusted according to the evaluation of the previous segment.

\begin{table}[t]
  \caption{Training setup.}
  \label{tab:training}
  \centering
  \small
  \begin{tabularx}{\linewidth}{@{}>{\raggedright\arraybackslash}p{4.4cm}>{\raggedright\arraybackslash}X@{}}
    \toprule
    Item & Setting \\
    \midrule
    Initialization & Qwen3-ASR-1.7B~\cite{shi2026qwen3asr} \\
    Stage 1: multi-task supervised fine-tuning & Plain ASR, hotword ASR, synthetic target-speaker mixtures and meeting turn-taking data; three rounds \\
    Stage 2: target-speaker robustness training & Five segments continued from the stage-1 model \\
    Parameters updated & All, except the text embeddings and the convolution layers of the AuT encoder, which are frozen; \septok is trained \\
    Loss & Cross-entropy on the text positions to be predicted only \\
    Released weights & Parameter-wise average of the last three checkpoints~\cite{wortsman2022model} \\
    \bottomrule
  \end{tabularx}
\end{table}

\paragraph{Weight averaging.}
The present-or-absent decision of a single checkpoint fluctuates: between neighboring checkpoints, the share of target-present samples with an empty output can differ by 4 percentage points.
We therefore release the parameter-wise average of the last three checkpoints~\cite{wortsman2022model}.
